% =============================================================================
% Bloque 7: Resultados (slides 17–24) — corazón de la defensa
% =============================================================================

\section{Resultados}

% --- Slide 17: HEADLINE ---
\begin{frame}{Resultado principal}
\begin{columns}[T]

\begin{column}{0.42\textwidth}
\vspace{0.4em}
\begin{headline}
\cifra{$+2.25$ pp}{vs.\ SMOTE\\ (macro F1)}
\end{headline}

\vspace{0.5em}
\small
\begin{itemize}\setlength\itemsep{2pt}
  \item $p < 0{.}0001$ \cita{Bonferroni}
  \item $d$ de Cohen $= 0{.}74$ \cita{efecto medio-grande}
  \item Win-rate $= 83{.}8\%$
  \item $N = 3{.}675$ configuraciones
  \item Clasificación absoluta: \textbf{73.49\%} vs.\ 71.24\%
\end{itemize}
\end{column}

\begin{column}{0.58\textwidth}
\centering
\includegraphics[width=\linewidth, height=6.5cm, keepaspectratio]{fig_boxplot_deltas.png}\\[2pt]
{\scriptsize Distribución de $\Delta$F1 por método.}
\end{column}

\end{columns}

\mensaje{El método propuesto supera todas las líneas base, incluidas las modernas.}
\end{frame}

% --- Slide 18a: Antes vs después del filtrado ---
\begin{frame}{Visualización t-SNE: el filtro en acción --- antes vs después}
\centering
\includegraphics[width=0.96\linewidth, height=5.4cm, keepaspectratio]{fig_tsne_before_after_20newsgroups.png}

\vspace{0.1em}
{\scriptsize \dataset{20newsgroups} (4 clases, 10-shot). \textbf{Izq.:} candidatas del LLM (triángulos tenues). \textbf{Der.:} solo las que pasaron el filtro. Círculos = reales, estrellas = centroides.}

\mensaje{El filtro elimina las candidatas en zonas de cruce o periféricas; conserva las que refuerzan el cluster.}
\end{frame}

% --- Slide 18b: t-SNE filtrado vs SMOTE ---
\begin{frame}{Visualización t-SNE: SMOTE vs filtro}
\centering
\includegraphics[width=0.85\linewidth, height=5.8cm, keepaspectratio]{fig_tsne_filtered_vs_smote_emotion.png}

\vspace{0.2em}
{\scriptsize Dataset \dataset{emotion} (6 clases, 10-shot): triángulos = LLM filtrado, cuadrados = SMOTE.}

\mensaje{El filtro rechaza candidatas en zonas de cruce y aporta diversidad real.}
\end{frame}

% --- Slide 19: F1 vs n-shot ---
\begin{frame}{Retornos decrecientes con más datos}
\begin{columns}[T]

\begin{column}{0.55\textwidth}
\centering
\includegraphics[width=\linewidth, height=5.8cm, keepaspectratio]{fig_delta_by_nshot.png}\\[2pt]
{\scriptsize $\Delta$F1 por valor de $n$-shot.}
\end{column}

\begin{column}{0.42\textwidth}
\small
\textbf{El método brilla donde más se necesita:}
\vspace{0.3em}

\begin{tabular}{@{}lcc@{}}
\toprule
$n$-shot & $\Delta$F1 & $d$ \\
\midrule
\textbf{10} & \alert{$+4{.}89$} & 1.48 \\
25          & $+2{.}17$         & 1.57 \\
50          & $+0{.}76$         & 0.80 \\
\bottomrule
\end{tabular}

\vspace{0.5em}
\begin{insight}
\scriptsize Cuanto menos datos hay, más útil es el filtrado.
\end{insight}
\end{column}

\end{columns}

\mensaje{Caso de uso ideal: 10 ejemplos por clase.}
\end{frame}

% --- Slide 20: Por clasificador ---
\begin{frame}{El clasificador importa: lineales ganan}
\begin{columns}[T]

\begin{column}{0.55\textwidth}
\centering
\includegraphics[width=\linewidth, height=5.5cm, keepaspectratio]{fig_classifier_comparison.png}
\end{column}

\begin{column}{0.42\textwidth}
\small
\begin{tabular}{@{}lr@{\hspace{6pt}}r@{}}
\toprule
Clasificador & $\Delta$F1 & $d$ \\
\midrule
SVC (lineal)        & \alert{$+2{.}99$} & 1.02 \\
Ridge               & $+2{.}69$         & 1.04 \\
MLP                 & $+2{.}16$         & 0.73 \\
Logistic Regression & $+1{.}98$         & 0.82 \\
Random Forest       & $+1{.}42$         & n.s. \\
\bottomrule
\end{tabular}

\vspace{0.5em}
{\scriptsize H1d confirmada: \textbf{ensambles de árboles}\\ no aprovechan los pesos suaves.}
\end{column}

\end{columns}

\mensaje{Modelos lineales explotan mejor la señal geométrica.}
\end{frame}

% --- Slide 21: Sweet spot por número de clases ---
\begin{frame}{¿Cuándo conviene aplicarlo? Por número de clases}
\begin{columns}[T]

\begin{column}{0.5\textwidth}
\centering
\includegraphics[width=\linewidth, height=5.5cm, keepaspectratio]{fig_heatmap_methods.png}
\end{column}

\begin{column}{0.5\textwidth}
\small
\textbf{Sweet spot: pocas clases.}
\vspace{0.2em}

\begin{tabular}{@{}lr@{}}
\toprule
\# Clases (dataset)     & $\Delta$F1 \\
\midrule
6 (\dataset{trec6, emotion}) & \alert{$+5{.}42$} \\
4 (\dataset{20news, ag})     & $+2{.}80$ \\
2 (binarios)                 & $+3{.}18$ \\
77 (\dataset{banking77})     & $-0{.}99$ \\
150 (\dataset{clinc150})     & $+0{.}01$ \\
\bottomrule
\end{tabular}

\vspace{0.4em}
{\scriptsize Spearman $\rho = -0{.}260$, $p=0{.}003$.}
\end{column}

\end{columns}

\mensaje{Saturación con $\geq 77$ clases: el espacio se fragmenta.}
\end{frame}

% --- Slide 22: Hallazgo contra-intuitivo ---
\begin{frame}{Hallazgo contra-intuitivo: simple $>$ complejo}
\begin{columns}[T]

\begin{column}{0.5\textwidth}
\centering
\begin{tikzpicture}[every node/.style={font=\scriptsize}, scale=0.7]
  % Línea vertical en 0 (referencia SMOTE)
  \draw[uaigray, dashed, thick] (0,-0.3) -- (0,4.8);
  \node[below, font=\tiny] at (0,-0.35) {SMOTE = 0};

  % Eje horizontal abajo
  \draw[->,thick,uaigray] (-3.7,-0.3) -- (2.2,-0.3) node[right, font=\tiny]{$\Delta$F1 (pp)};
  \foreach \x in {-3,-2,-1,1,2}{
    \draw[uaigray!50] (\x,-0.22) -- (\x,-0.38);
  }

  % Bars (positive)
  \fill[methodOurs] (0,4.1) rectangle (1.45,4.5);
  \node[right, font=\tiny] at (1.55,4.3) {\textbf{Cascade L1 +1.45}};

  \fill[methodOurs!70] (0,3.4) rectangle (0.88,3.8);
  \node[right, font=\tiny] at (0.98,3.6) {LOF +0.88};

  \fill[methodOther] (0,2.7) rectangle (0.72,3.1);
  \node[right, font=\tiny] at (0.82,2.9) {None +0.72};

  % Bars (negative)
  \fill[methodWorse] (-1.29,2.0) rectangle (0,2.4);
  \node[left, font=\tiny] at (-1.39,2.2) {Emb.-Guided $-1.29$};

  \fill[methodWorse] (-3.28,1.3) rectangle (0,1.7);
  \node[left, font=\tiny, color=uaiorange] at (-3.38,1.5) {\textbf{Combined $-3.28$}};
\end{tikzpicture}
\end{column}

\begin{column}{0.5\textwidth}
\setbeamercolor{block title}{fg=white,bg=uaiamber}
\begin{block}{Insight clave}
Filtros multi-criterio \textbf{degradan} el rendimiento.
\vspace{0.4em}
\begin{itemize}\setlength\itemsep{0pt}\scriptsize
  \item Distancia y similitud son redundantes en L2-norm.
  \item LOF inestable con $n=10$.
  \item Restricciones conjuntivas matan la diversidad.
\end{itemize}
\end{block}

\vspace{0.4em}
\begin{insight}
\scriptsize\centering
Más filtrado $\neq$ mejor filtrado.\\Generar muchos, filtrar poco.
\end{insight}
\end{column}

\end{columns}
\end{frame}

% --- Slide 23: Robustez multi-LLM y multi-embedding ---
\begin{frame}{Robustez: ¿depende del LLM o del embedding?}
\begin{columns}[T]

\begin{column}{0.5\textwidth}
\centering
\includegraphics[width=\linewidth, height=4.6cm, keepaspectratio]{fig_multi_llm_comparison.png}\\[2pt]
\textbf{5 LLMs probados}\\
{\scriptsize Gemini, GPT-5, Claude, Kimi, GLM}\\
{\scriptsize $\Delta$F1 entre $+1{.}49$ y $+2{.}31$ pp}\\
{\scriptsize Friedman $p=0{.}086$ \cita{no significativo}}
\end{column}

\begin{column}{0.5\textwidth}
\centering
\setbeamercolor{block title}{fg=white,bg=uaiblue}
\begin{block}{4 modelos de embeddings}
\scriptsize
\begin{tabular}{@{}lr@{}}
\toprule
Modelo & $\Delta$F1 \\
\midrule
mpnet-base-v2 (768d)     & $+2{.}74$ \\
bge-large (1024d)        & $+2{.}04$ \\
bge-small (384d)         & $+2{.}00$ \\
e5-large (1024d)         & $+0{.}97$ \\
\bottomrule
\end{tabular}
\vspace{0.3em}\\
{\scriptsize 90.5\% acuerdo entre modelos.}
\end{block}
\end{column}

\end{columns}

\mensaje{El método es agnóstico al LLM y al modelo de embeddings.}
\end{frame}

% --- Slide 24: NER ---
\begin{frame}{Generalización a NER: \textit{el filtro funciona en otra tarea}}
\begin{columns}[T]

\begin{column}{0.45\textwidth}
\small
\textbf{Reconocimiento de entidades nombradas:}
\begin{itemize}\setlength\itemsep{0pt}
  \item 3 corpus: \dataset{MultiNERD}, \dataset{WikiANN}, \dataset{Few-NERD}
  \item Adaptación: clase = tipo de entidad dominante por oración.
  \item \textbf{Filtros sin modificación.}
\end{itemize}

\vspace{0.4em}
\begin{headline}
\cifra{$+9{.}61$ pp}{Cascade L1\\ ($p=0.0003$, $d=2.05$)}
\end{headline}
\end{column}

\begin{column}{0.55\textwidth}
\centering
\begin{tikzpicture}[every node/.style={font=\scriptsize}, x=0.32cm, y=0.5cm]
  % Eje horizontal y vertical
  \draw[->,thick,uaigray] (0,0) -- (16.5,0) node[right, font=\tiny]{$\Delta$F1 (pp)};
  \draw[uaigray] (0,0) -- (0,7.5);
  \foreach \x in {0,3,6,9,12,15}{
    \draw[uaigray!50] (\x,0.15) -- (\x,-0.15) node[below, font=\tiny]{\x};
  }

  % Cascade L1
  \fill[methodOurs] (0,6) rectangle (9.61,7);
  \node[right, font=\scriptsize\bfseries] at (9.81,6.5) {Cascade L1 +9.61};

  % LOF
  \fill[methodOurs!80] (0,4.5) rectangle (9.28,5.5);
  \node[right, font=\scriptsize] at (9.48,5.0) {LOF +9.28};

  % None: solo línea de referencia (sin barra)
  \draw[methodSMOTE, very thick] (-0.3,3) -- (0.3,4);
  \node[right, font=\scriptsize, color=methodSMOTE] at (0.5,3.5) {None (referencia $= 0$)};

  % Combined: marca de rechazo (no medido)
  \node[right, font=\scriptsize, color=uaiorange] at (0.2,2) {\textbf{$\times$} Combined: 0\% aceptación};

  % Anotación 10-shot
  \node[font=\tiny, color=uaigreen!50!black] at (8,0.8) {(10-shot: $+15.15$ pp)};
\end{tikzpicture}
\end{column}

\end{columns}

\mensaje{La calidad geométrica es agnóstica a la tarea.}
\end{frame}
